\section{ĐÁNH GIÁ HIỆU QUẢ}

\subsection{Tiêu chí đánh giá}

Hiệu quả của hệ thống được đánh giá theo hai nhóm tiêu chí. Nhóm thứ nhất là độ chính xác chẩn đoán: detection score đo khả năng phát hiện có lỗi hay không; localization accuracy/F1 đo khả năng xác định đúng thiết bị lỗi; RCA accuracy/F1 đo khả năng xác định đúng nguyên nhân gốc rễ; incident success rate yêu cầu đúng đồng thời cả detection, localization và RCA. Nhóm thứ hai là hiệu năng vận hành: số bước suy luận, số lần gọi công cụ và lượng token tiêu thụ trung bình mỗi incident.

Ba tầng metric này được chọn vì phản ánh quy trình xử lý sự cố trong thực tế. Detection đúng nhưng localization sai vẫn chưa đủ để kỹ sư xử lý sự cố; localization đúng nhưng RCA sai có thể dẫn đến thao tác khắc phục không hiệu quả. Vì vậy, incident success rate là tiêu chí khắt khe nhất, còn F1 ở localization/RCA giúp đánh giá các trường hợp agent tìm đúng một phần tập thiết bị hoặc nhãn nguyên nhân. Việc đo thêm steps, tool calls và tokens giúp phân biệt giữa cải thiện do suy luận tốt hơn và cải thiện phải trả bằng chi phí vận hành quá lớn.

Hai cấu hình được so sánh trên cùng tập incident:
\begin{itemize}
  \item \textbf{Baseline Agent}: ReAct agent tiêu chuẩn, không dùng memory và không dùng tool evolution.
  \item \textbf{Evolved Agent}: ReAct agent tích hợp Procedural Memory và Tool Refinement, có khả năng học liên tục qua các episode.
\end{itemize}

\subsection{Kết quả so sánh}

\begin{table*}[t]
\centering
\small
\caption{Kết quả chung của ReAct + GPT-OSS-120B trên 125 case chung}
\begin{tabularx}{\textwidth}{@{}Yccc@{}}
\toprule
\textbf{Chỉ số} & \textbf{Baseline} & \textbf{Evolved} & \textbf{Thay đổi} \\
\midrule
Detection score & 0.840 & 0.872 & +3.8\% \\
Localization accuracy & 0.368 & 0.384 & +4.3\% \\
Localization F1 & 0.411 & 0.459 & +11.8\% \\
RCA accuracy & 0.304 & 0.376 & +23.7\% \\
RCA F1 & 0.304 & 0.387 & +27.2\% \\
Incident success rate & 26/125 (20.8\%) & 35/125 (28.0\%) & +34.6\% \\
\midrule
Steps trung bình & 15.29 & 12.74 & -16.7\% \\
Tool calls trung bình & 13.60 & 11.09 & -18.5\% \\
Tool error rate & 11/1700 (0.65\%) & 0/1386 (0.00\%) & -100.0\% \\
Input tokens (tổng) & 15.944M & 17.992M & +12.9\% \\
Output tokens (tổng) & 0.697M & 0.708M & +1.5\% \\
Total tokens (tổng) & 16.641M & 18.699M & +12.4\% \\
Tổng thời gian chạy & 207.7 phút & 222.6 phút & +7.2\% \\
\bottomrule
\end{tabularx}
\end{table*}

Kết quả cho thấy agent tự tiến hóa cải thiện rõ nhất ở RCA và incident success rate: RCA F1 tăng từ 0.304 lên 0.387, số incident giải quyết đúng hoàn toàn tăng từ 26/125 lên 35/125. Evolved Agent cũng giảm số bước 16.7\%, giảm tool calls 18.5\% và đưa tool error rate về 0.00\%. Đổi lại, tổng token tăng 12.4\% và tổng thời gian chạy tăng 7.2\%, chủ yếu do prompt phải nạp thêm kỹ năng và tài liệu công cụ đã tinh chỉnh.

\begin{table}[t]
\centering
\scriptsize
\caption{Placeholder: xu hướng RCA/Localization theo số case tăng dần. Số liệu minh họa, chưa dùng làm kết quả chính thức.}
\label{tab:placeholder_case_curve}
\begin{tabularx}{\columnwidth}{@{}c c c Y@{}}
\toprule
\textbf{Phạm vi case} & \textbf{RCA F1} & \textbf{Loc. F1} & \textbf{Diễn giải}\\
\midrule
1--25 & .43 & .55 & Tăng nhanh nhờ nhiều case liên quan.\\
1--50 & .42 & .53 & Hiệu quả còn rõ nhưng tốc độ tăng chậm hơn.\\
1--75 & .40 & .49 & Bắt đầu gặp topology và tín hiệu lỗi xa hơn.\\
1--125 & .39 & .46 & Metric hội tụ thấp hơn so với giai đoạn đầu.\\
\bottomrule
\end{tabularx}
\vspace{0.15em}
\begin{minipage}{\columnwidth}
\footnotesize\textit{Phân tích.} Metric tăng mạnh ở giai đoạn đầu vì Procedural Memory tái sử dụng được quy trình từ các case liên quan; lợi ích giảm khi gặp topology hoặc nhóm lỗi xa hơn.
\end{minipage}
\end{table}

\begin{table}[t]
\centering
\scriptsize
\caption{Placeholder: độ ổn định của ReAct + GPT-OSS-120B qua 3 lần chạy. Số liệu minh họa, chưa dùng làm kết quả chính thức.}
\label{tab:placeholder_stability}
\begin{tabularx}{\columnwidth}{@{}c cc cc Y@{}}
\toprule
\textbf{Run} & \multicolumn{2}{c}{\textbf{Success rate}} & \multicolumn{2}{c}{\textbf{RCA F1}} & \textbf{Nhận xét}\\
\cmidrule(lr){2-3}\cmidrule(lr){4-5}
 & \textbf{Base} & \textbf{Evol.} & \textbf{Base} & \textbf{Evol.} & \\
\midrule
1 & 20.8 & 28.0 & .304 & .387 & Run chính thức.\\
2 & 19.6 & 27.2 & .296 & .381 & Khoảng cách vẫn dương.\\
3 & 21.6 & 29.6 & .311 & .395 & Không đảo chiều kết luận.\\
\bottomrule
\end{tabularx}
\vspace{0.15em}
\begin{minipage}{\columnwidth}
\footnotesize\textit{Phân tích.} Evolved Agent được kỳ vọng ổn định hơn baseline nếu độ lệch giữa các run nhỏ nhưng vẫn giữ khoảng cách hiệu năng dương; bảng chính thức sẽ thay bằng trung bình và độ lệch chuẩn.
\end{minipage}
\end{table}

\begin{table}[t]
\centering
\scriptsize
\caption{Placeholder: so sánh chiến lược agent trên GPT-OSS-120B. Số liệu minh họa, chưa dùng làm kết quả chính thức.}
\label{tab:placeholder_strategy}
\begin{tabularx}{\columnwidth}{@{}l c c c Y@{}}
\toprule
\textbf{Agent} & \textbf{Success} & \textbf{RCA F1} & \textbf{Calls} & \textbf{Diễn giải}\\
\midrule
ReAct & 28.0\% & .387 & 11.09 & Vòng lặp quan sát--hành động ngắn.\\
Plan\&Execute & 25.6\% & .361 & 12.74 & Có cấu trúc nhưng dễ lệch khi quan sát thay đổi.\\
Reflexion & 26.4\% & .374 & 13.28 & Phản tư thêm nhưng tốn bước và token hơn.\\
\bottomrule
\end{tabularx}
\vspace{0.15em}
\begin{minipage}{\columnwidth}
\footnotesize\textit{Phân tích.} Placeholder này kiểm tra luận điểm ReAct phù hợp hơn với benchmark hiện tại do có vòng lặp quan sát--hành động ngắn; Plan\&Execute và Reflexion có thêm chi phí lập kế hoạch hoặc phản tư.
\end{minipage}
\end{table}

\begin{table}[t]
\centering
\scriptsize
\caption{Placeholder: so sánh ReAct trên GPT-OSS-120B và Qwen3.5-122B-A10B-FP8. Số liệu minh họa, chưa dùng làm kết quả chính thức.}
\label{tab:placeholder_model}
\begin{tabularx}{\columnwidth}{@{}l cc cc Y@{}}
\toprule
\textbf{Model} & \multicolumn{2}{c}{\textbf{Baseline}} & \multicolumn{2}{c}{\textbf{Evolved}} & \textbf{Nhận xét}\\
\cmidrule(lr){2-3}\cmidrule(lr){4-5}
 & \textbf{Success} & \textbf{RCA} & \textbf{Success} & \textbf{RCA} & \\
\midrule
GPT-OSS & 20.8 & .304 & 28.0 & .387 & Model chính.\\
Qwen3.5 & 18.4 & .286 & 25.6 & .365 & Backbone đối chứng.\\
\bottomrule
\end{tabularx}
\vspace{0.15em}
\begin{minipage}{\columnwidth}
\footnotesize\textit{Phân tích.} So sánh model giúp tách ảnh hưởng của mô hình nền khỏi lợi ích của Procedural Memory và Tool Refinement; nếu cả hai model đều tăng khi bật evolved mode, đóng góp hệ thống thuyết phục hơn.
\end{minipage}
\end{table}

\subsection{Phân tích theo nhóm lỗi}

Evolved Agent cải thiện mạnh nhất ở nhóm lỗi host như \code{host\_missing\_ip}, \code{host\_incorrect\_ip}, \code{host\_incorrect\_gateway}. Các kỹ năng tích lũy từ case trước giúp agent nhanh chóng kiểm tra cấu hình host, subnet, gateway và reachability thay vì thử nhiều hướng không liên quan. Với các lỗi định tuyến BGP/OSPF phức tạp, hệ thống có cải thiện nhưng vẫn chưa ổn định vì cần suy luận đồng thời trên neighbor state, route advertisement và quan hệ topology. Nhóm resource contention và một số lỗi tấn công vẫn khó do tín hiệu quan sát không trực tiếp, cần thêm telemetry hoặc cơ chế kiểm chứng giả thuyết mạnh hơn.

Nhìn chung, hai mô-đun tiến hóa tạo ra hiệu quả bổ trợ: Procedural Memory cải thiện chiến lược chẩn đoán, Tool Refinement cải thiện cách dùng công cụ. Đây là điểm khác biệt so với agent tĩnh chỉ dựa vào prompt ban đầu và lịch sử hội thoại ngắn hạn.

\subsection{Hạn chế và rủi ro}

Hệ thống vẫn có một số hạn chế. Thứ nhất, hiệu quả tự tiến hóa phụ thuộc vào chất lượng episode trước đó; nếu feedback hoặc ground truth không đáng tin cậy, bộ nhớ có thể ghi nhận quy trình sai. Thứ hai, các kỹ năng học từ một nhóm topology có thể quá khớp với tên thiết bị, giao thức hoặc cách tổ chức lab cũ, làm giảm khả năng chuyển sang topology mới. Thứ ba, Tool Refinement cải thiện mô tả công cụ nhưng không làm tăng năng lực quan sát nếu môi trường không cung cấp telemetry hoặc log cần thiết. Cuối cùng, việc đưa thêm kỹ năng và tài liệu vào prompt làm tăng token, nên cần cơ chế chọn lọc ngữ cảnh chặt chẽ hơn khi triển khai ở quy mô lớn.

Các rủi ro này cho thấy self-evolving agent cần được vận hành theo nguyên tắc có kiểm soát. Trong môi trường mạng quan trọng, hệ thống nên ghi rõ bằng chứng cho từng kết luận, cho phép kỹ sư phê duyệt kỹ năng mới trước khi lưu lâu dài và có cơ chế rollback bộ nhớ khi phát hiện tri thức sai. Đây cũng là lý do đề tài ưu tiên tiến hóa bộ nhớ/tài liệu thay vì tự động sinh thao tác cấu hình mới trên thiết bị mạng.
